%%%%%%%%%%%%%%%%%
% CV tailored for Thales - WPM Management de Work Packages et support aux bids/captures
%%%%%%%%%%%%%%%%%

\documentclass[8pt,a4paper,ragged2e,withhyper]{altacv}

\geometry{left=1.25cm,right=1.25cm,top=.5cm,bottom=1.cm,columnsep=1.2cm}

\usepackage{paracol}
\usepackage{fontawesome5}
\usepackage{hyperref}

\iftutex 
  \setmainfont{Roboto Slab}
  \setsansfont{Lato}
  \renewcommand{\familydefault}{\sfdefault}
\else
  \usepackage[rm]{roboto}
  \usepackage[defaultsans]{lato}
  \renewcommand{\familydefault}{\sfdefault}
\fi

\definecolor{SlateGrey}{HTML}{2E2E2E}
\definecolor{LightGrey}{HTML}{666666}
\definecolor{DarkPastelRed}{HTML}{8AC0D9}
\definecolor{PastelRed}{HTML}{00B0DA}
\definecolor{GoldenEarth}{HTML}{B4AD0C}

\colorlet{name}{black}
\colorlet{tagline}{PastelRed}
\colorlet{heading}{DarkPastelRed}
\colorlet{headingrule}{GoldenEarth}
\colorlet{subheading}{PastelRed}
\colorlet{accent}{PastelRed}
\colorlet{emphasis}{SlateGrey}
\colorlet{body}{LightGrey}

\definecolor{violet}{HTML}{5A2582}
\definecolor{rouge}{HTML}{F53B3B}
\definecolor{noir}{HTML}{00232C}
\definecolor{bleu}{HTML}{00B0DA}
\definecolor{rose}{HTML}{F831A0}
\definecolor{jaune}{HTML}{FFEC00}
\definecolor{vert}{HTML}{34AD6C}

\renewcommand{\namefont}{\Huge\rmfamily\bfseries}
\renewcommand{\personalinfofont}{\footnotesize}
\renewcommand{\cvsectionfont}{\LARGE\rmfamily\bfseries}
\renewcommand{\cvsubsectionfont}{\large\bfseries}
\renewcommand{\cvItemMarker}{{\small\textbullet}}
\renewcommand{\cvRatingMarker}{\faCircle}

\input{pubs-num.tex}
\addbibresource{sample.bib}

\begin{document}

\name{BOILEAU Guillaume}
\tagline{WPM / Ingénieur coordination technique | Spatial, IVVQ, performance système \& interfaces partenaires}

\photoL{2.8cm}{moi}

\personalinfo{%
  \email{guillaume.boileaupro@gmail.com}
  \phone{+33 6 59 94 85 84}
  \location{Alpes-Maritimes, FRANCE}
  \linkedin{guillaume-boileau-891b81265}
  \researchgate{Guillaume-Boileau-2}
  \orcid{0000-0002-3576-6968}
}

\makecvheader
\vspace{-0.3cm}

\AtBeginEnvironment{itemize}{\small}
\columnratio{0.64}

\begin{paracol}{2}

\cvsection{Experience}

\cvevent{\textbf{Ingénieur logiciel full stack -- Intégration \& support opérationnel}}{VEV Platform Services France}{2025 -- 2026}{Valbonne, France}

\textbf{Contexte :} Développement logiciel et support opérationnel pour une plateforme CPMS liée à des infrastructures de recharge de véhicules électriques.

\textbf{Objectif :} Contribuer au développement, à l’intégration et au suivi technique de services logiciels connectés à des équipements terrain et à des flux de données opérationnels.

\begin{itemize}
  \item \textbf{Développement logiciel :} Contribution au développement d’applications web et de services backend en TypeScript, Angular et Node.js.
  \item \textbf{Intégration système :} Travail sur les interfaces entre services logiciels, données opérationnelles et infrastructures connectées.
  \item \textbf{Suivi opérationnel :} Analyse de flux de données, notamment via des échanges FTP utilisés pour le pilotage et le suivi technique.
  \item \textbf{Support production :} Analyse d’incidents, investigation de comportements applicatifs et contribution à la fiabilisation des services.
  \item \textbf{Interfaces terrain :} Vérification de la cohérence entre données applicatives, comportement des équipements et attentes opérationnelles.
  \item \textbf{Coordination agile :} Utilisation de Zenhub pour le suivi des tâches, la priorisation, le backlog et la coordination dans un environnement itératif.
  \item \textbf{Reporting technique :} Formalisation d’observations, d’analyses d’incidents et de points techniques pour faciliter la prise de décision.
\end{itemize}

\divider

\cvevent{\textbf{Ingénieur R\&D senior -- Coordination technique, performance \& validation}}{Laboratoire Lagrange, Observatoire de la Côte d’Azur, CNES}{2023 -- 2025}{Nice, France}

\textbf{Contexte :} Activités d’ingénierie et de R\&D pour la mission spatiale LISA ESA/NASA, dans un environnement international impliquant simulation, traitement de données, intégration de modèles et validation technique.

\textbf{Objectif :} Développer et structurer des outils Python permettant d’intégrer, comparer, valider et documenter des contributions techniques hétérogènes.

\begin{itemize}
  \item \textbf{Coordination technique :} Interface avec des contributeurs scientifiques, logiciels et institutionnels dans un environnement spatial international.
  \item \textbf{Intégration de contributions :} Consolidation de résultats issus de modèles différents dans un cadre commun d’analyse, de comparaison et de revue.
  \item \textbf{Performance système :} Analyse de résultats de simulation, comparaison de modèles et évaluation de cohérence dans des contextes de signaux faibles et de bruit.
  \item \textbf{IVVQ / validation :} Définition de contrôles de cohérence, de règles de comparaison et de méthodes de vérification des résultats.
  \item \textbf{Plan de développement logiciel :} Structuration de workflows, versions, configurations, sorties de simulation et résultats d’analyse.
  \item \textbf{Documentation technique :} Utilisation de Confluence pour organiser la documentation, la traçabilité des analyses, le suivi de validation et le partage de connaissances.
  \item \textbf{Support aux revues :} Production de synthèses techniques, résultats traçables et éléments de comparaison pour revue collaborative.
  \item \textbf{Plateforme logicielle :} Développement de CATpyLISA, plateforme Python pour l’intégration, la comparaison, la validation et la revue technique de modèles.
  \textcolor{DarkPastelRed}{\href{https://gitlab.oca.eu/gboileau/lisa_l3_tools}{\bf GitLab:CATpyLISA}}
\end{itemize}

\divider

\cvevent{\textbf{Data Scientist -- Modélisation, bruit \& analyse de performance}}{Universiteit Antwerpen, Belgium}{2021 -- 2023}{Antwerp, Belgium}

\textbf{Contexte :} Analyse de données et études de performance pour des instruments de précision dans un environnement technique international.

\textbf{Objectif :} Développer des workflows robustes et reproductibles pour identifier, caractériser et réduire des sources affectant la qualité de mesure.

\begin{itemize}
  \item \textbf{Analyse de performance :} Évaluation de la qualité de mesure, des limites instrumentales et de l’impact de différentes sources de bruit.
  \item \textbf{Traitement du signal :} Identification et caractérisation de contributions de bruit dans des mesures complexes.
  \item \textbf{Modélisation statistique :} Développement de méthodes d’analyse et d’inférence bayésienne, notamment par chaînes MCMC.
  \item \textbf{Pipelines Python :} Mise en place de workflows de traitement, filtrage, transformation et analyse statistique.
  \item \textbf{Validation de résultats :} Analyse de sensibilité, vérification de cohérence et évaluation des limites de performance.
  \item \textbf{Reporting technique :} Production de résultats traçables, de synthèses techniques et de supports pour des parties prenantes internationales.
\end{itemize}

\divider

\cvevent{\textbf{Ingénieur R\&D junior -- Signal, simulation \& validation}}{ARTEMIS, Observatoire de la Côte d’Azur}{2018 -- 2021}{Nice, France}

\textbf{Contexte :} Activités de R\&D liées à l’analyse de signaux faibles, à la simulation numérique et à la préparation de missions spatiales et d’instruments de détection.

\textbf{Objectif :} Développer des méthodes d’analyse, automatiser des simulations et produire des résultats exploitables dans un environnement scientifique et technique exigeant.

\begin{itemize}
  \item \textbf{Simulation numérique :} Développement d’approches de simulation pour des environnements de signaux complexes et des jeux de données de grande taille.
  \item \textbf{Analyse de données :} Mise en place de méthodes pour analyser et séparer des signaux faibles et superposés.
  \item \textbf{Traitement du signal :} Études de caractérisation, de décomposition et de comparaison de signaux dans des contextes bruités.
  \item \textbf{Validation technique :} Vérification de cohérence, comparaison de résultats et études de sensibilité.
  \item \textbf{Investigation technique :} Analyse de sources de bruit et réalisation de diagnostics par corrélations entre capteurs.
  \item \textbf{Documentation :} Formalisation des méthodes, hypothèses, résultats et conclusions pour revue collaborative.
\end{itemize}

\switchcolumn

\cvsection{Profile}

\textbf{Ingénieur docteur avec une expérience solide en environnement spatial, coordination technique, validation de résultats, performance système, traitement du signal et développement d’outils logiciels pour systèmes complexes.}

Mon parcours combine simulation numérique, analyse de performance, intégration de contributions hétérogènes, validation technique, documentation, coordination avec parties prenantes et support à la décision. J’ai travaillé dans des environnements exigeants : mission spatiale LISA ESA/NASA, CNES, collaborations internationales, instrumentation scientifique et plateformes logicielles industrielles.

\medskip

\textbf{Positionnement cible :} Work Package Manager, ingénieur coordination technique, ingénieur IVVQ système, responsable de lots techniques, support bids/captures ou coordination R\&T dans le spatial.

\begin{itemize}
  \item Capacité à piloter des activités transverses entre données, logiciel, modèles, exigences, interfaces et parties prenantes.
  \item Expérience en validation technique, contrôles de cohérence, simulation, analyse de performance et traçabilité des résultats.
  \item Culture forte du spatial, des environnements institutionnels, des revues techniques et de la collaboration internationale.
  \item Profil rigoureux, calme, orienté livraison, documentation, analyse des risques et respect des jalons.
\end{itemize}

\cvsection{Languages}

\cvskill{English}{4}
\divider
\cvskill{Spanish}{2}
\divider

\medskip

\cvsection{Education}

\cvevent{PhD in Physics}{Université Côte d’Azur}{2018 -- 2021}{Nice, France}

\divider

\cvevent{Selective Graduate Program in Fundamental Physics (Magistère)}{Université Paris-Saclay}{2016 -- 2018}{Orsay, France}

\divider

\cvevent{MSc in Astronomy and Astrophysics}{Observatoire de Paris / Université Paris-Saclay}{2017 -- 2018}{Paris, France}

\divider

\cvevent{BSc in Fundamental Physics}{Université Paris-Saclay}{2015 -- 2016}{Orsay, France}

\cvsection{Professional Training}

\cvevent{Supervised Machine Learning: Regression \& Classification}{DeepLearning.AI -- Andrew Ng}{2026}{Coursera}
\divider

\cvevent{Développement assisté par IA}{Prompting, relecture de code, debugging, prompting incrémental et critique}{2026}{Autoformation}

\end{paracol}

\cvsection{Projects}

\cvevent{CATpyLISA -- Plateforme Python d’intégration, comparaison et validation}{Mission spatiale LISA ESA/NASA}{2023 -- 2025}{}

Développement d’une plateforme Python dédiée à l’intégration de modèles, à la structuration de données, à la comparaison de résultats, à la validation technique et à la revue de workflows liés à la mission LISA.

\textbf{Rôle :} développement Python, conception de pipelines, structuration de données, logique de validation, contrôles de cohérence, documentation technique et coordination avec plusieurs contributeurs.

\textbf{Compétences mobilisées :} Python, intégration de modèles, validation, traitement du signal, analyse de performance, simulation, documentation, GitLab, environnement spatial international.

\divider

\cvevent{Workflows de simulation, performance et analyse de bruit}{Instrumentation scientifique et préparation de missions spatiales}{2018 -- 2025}{}

Développement de workflows numériques et statistiques pour analyser des signaux faibles, caractériser des contributions de bruit, comparer des résultats et évaluer la robustesse de modèles dans des environnements complexes.

\textbf{Rôle :} analyse de performance, simulation, traitement du signal, caractérisation de bruit, études de sensibilité, diagnostics techniques et reporting.

\textbf{Compétences mobilisées :} traitement du signal, simulation numérique, analyse spectrale, Python, C, validation, analyse de sensibilité, documentation technique.

\cvsection{Compétences clés}

\vspace{-.3cm}

\cvsubsection{Work Package Management, pilotage \& coordination}

{\small
\cvtag{Work Package Management}
\cvtag{Coordination technique}
\cvtag{Gestion de lots techniques}
\cvtag{Pilotage transverse}
\cvtag{Interfaces clients}
\cvtag{Interfaces partenaires}
\cvtag{Interfaces fournisseurs}
\cvtag{Suivi de jalons}
\cvtag{Reporting}
\cvtag{KPIs}
\cvtag{Planning}
\cvtag{Risques \& opportunités}
\cvtag{Support décisionnel}
\cvtag{REX}
\cvtag{Amélioration continue}
}

\divider\smallskip
\vspace{-.3cm}

\cvsubsection{Bids, offres, conformité \& documentation}

{\small
\cvtag{Support RAO}
\cvtag{Support bids/captures}
\cvtag{Documentation technique}
\cvtag{Revues techniques}
\cvtag{Matrice de conformité}
\cvtag{Traçabilité}
\cvtag{Analyse macroscopique}
\cvtag{Justification technique}
\cvtag{Synthèse technique}
\cvtag{Confluence}
\cvtag{Jira}
\cvtag{Zenhub}
\cvtag{Anglais technique}
}

\divider\smallskip
\vspace{-.3cm}

\cvsubsection{IVVQ, validation, qualité \& systèmes}

{\small
\cvtag{IVVQ système}
\cvtag{IVVQ logiciel}
\cvtag{Vérification / Validation}
\cvtag{Plan de développement}
\cvtag{Cycle en V}
\cvtag{Agile}
\cvtag{Tests fonctionnels}
\cvtag{Non-régression}
\cvtag{Contrôles de cohérence}
\cvtag{Analyse d’anomalies}
\cvtag{Root cause analysis}
\cvtag{Diagnostic technique}
\cvtag{Documentation qualité}
\cvtag{Systèmes complexes}
}

\divider\smallskip
\vspace{-.3cm}

\cvsubsection{Performance système, simulation \& signal}

{\small
\cvtag{Performance système}
\cvtag{Analyse de performance}
\cvtag{Simulation numérique}
\cvtag{Modélisation}
\cvtag{Traitement du signal}
\cvtag{Analyse spectrale}
\cvtag{Analyse de bruit}
\cvtag{Signaux faibles}
\cvtag{Low-SNR}
\cvtag{Analyse de sensibilité}
\cvtag{Comparaison de modèles}
\cvtag{Data quality}
\cvtag{Validation de résultats}
\cvtag{Calibration awareness}
\cvtag{Chaîne image awareness}
}

\divider\smallskip
\vspace{-.3cm}

\cvsubsection{Logiciel, data \& automatisation}

{\small
\cvtag{Python}
\cvtag{C}
\cvtag{Pandas}
\cvtag{NumPy}
\cvtag{SciPy}
\cvtag{SQL}
\cvtag{TypeScript}
\cvtag{Angular}
\cvtag{Node.js}
\cvtag{REST API}
\cvtag{Bash}
\cvtag{Linux}
\cvtag{Git}
\cvtag{GitLab}
\cvtag{Docker}
\cvtag{CI/CD}
\cvtag{FTP}
\cvtag{Power BI}
}

\divider\smallskip
\vspace{-.3cm}

\cvsubsection{Domaines d’application}

{\small
\cvtag{Spatial}
\cvtag{Systèmes satellitaires}
\cvtag{LISA ESA/NASA}
\cvtag{CNES}
\cvtag{ESA}
\cvtag{Instrumentation scientifique}
\cvtag{Observation de la Terre awareness}
\cvtag{Chaîne bord-sol awareness}
\cvtag{Systèmes complexes}
\cvtag{Projets internationaux}
}

\end{document}